\subsection{ソフトウェアレビューと評価}
レビューと評価は、定められた時点または必要な時点で、プロジェクトが目的へ進んでいるか、利害関係者の要求を満たしているか、プロセス、人員、ツール、方法が有効かを確認する活動である。

\subsubsection{要求充足の判定}
利害関係者の満足を達成することは、ソフトウェアエンジニアリング管理者の主要な目標である。進捗は、主要マイルストーン、ソフトウェア技術レビュー、反復サイクルの完了、製品増分の完成時に評価する。

要求からの差異を見つけ、適切な対応を取る。対応には、要求見直し、設計変更、追加テスト、受け入れ基準の明確化、範囲調整が含まれる。構成管理手順に従い、判断と変更の理由を記録する。

\subsubsection{パフォーマンスのレビューと評価}
プロジェクト人員の定期的なパフォーマンスレビューは、計画やプロセスに従える可能性、チーム内の困難、対立、支援の必要性を把握するために役立つ。

プロジェクトで用いる方法、ツール、技法も、有効性と適切性の観点で評価する。プロセスは、状況に対して意味があるか、役に立つか、効果があるかを定期的に確認し、必要に応じて変更する。
